\documentclass[
  man,
  floatsintext,
  longtable,
  nolmodern,
  notxfonts,
  notimes,
  colorlinks=true,linkcolor=blue,citecolor=blue,urlcolor=blue]{apa7}

\usepackage{amsmath}
\usepackage{amssymb}




\RequirePackage{longtable}
\RequirePackage{threeparttablex}

\makeatletter
\renewcommand{\paragraph}{\@startsection{paragraph}{4}{\parindent}%
	{0\baselineskip \@plus 0.2ex \@minus 0.2ex}%
	{-.5em}%
	{\normalfont\normalsize\bfseries\typesectitle}}

\renewcommand{\subparagraph}[1]{\@startsection{subparagraph}{5}{0.5em}%
	{0\baselineskip \@plus 0.2ex \@minus 0.2ex}%
	{-\z@\relax}%
	{\normalfont\normalsize\bfseries\itshape\hspace{\parindent}{#1}\textit{\addperi}}{\relax}}
\makeatother




\usepackage{longtable, booktabs, multirow, multicol, colortbl, hhline, caption, array, float, xpatch}
\setcounter{topnumber}{2}
\setcounter{bottomnumber}{2}
\setcounter{totalnumber}{4}
\renewcommand{\topfraction}{0.85}
\renewcommand{\bottomfraction}{0.85}
\renewcommand{\textfraction}{0.15}
\renewcommand{\floatpagefraction}{0.7}

\usepackage{tcolorbox}
\tcbuselibrary{listings,theorems, breakable, skins}
\usepackage{fontawesome5}

\definecolor{quarto-callout-color}{HTML}{909090}
\definecolor{quarto-callout-note-color}{HTML}{0758E5}
\definecolor{quarto-callout-important-color}{HTML}{CC1914}
\definecolor{quarto-callout-warning-color}{HTML}{EB9113}
\definecolor{quarto-callout-tip-color}{HTML}{00A047}
\definecolor{quarto-callout-caution-color}{HTML}{FC5300}
\definecolor{quarto-callout-color-frame}{HTML}{ACACAC}
\definecolor{quarto-callout-note-color-frame}{HTML}{4582EC}
\definecolor{quarto-callout-important-color-frame}{HTML}{D9534F}
\definecolor{quarto-callout-warning-color-frame}{HTML}{F0AD4E}
\definecolor{quarto-callout-tip-color-frame}{HTML}{02B875}
\definecolor{quarto-callout-caution-color-frame}{HTML}{FD7E14}

%\newlength\Oldarrayrulewidth
%\newlength\Oldtabcolsep


\usepackage{hyperref}




\providecommand{\tightlist}{%
  \setlength{\itemsep}{0pt}\setlength{\parskip}{0pt}}
\usepackage{longtable,booktabs,array}
\newcounter{none} % for unnumbered tables
\usepackage{calc} % for calculating minipage widths
% Correct order of tables after \paragraph or \subparagraph
\usepackage{etoolbox}
\makeatletter
\patchcmd\longtable{\par}{\if@noskipsec\mbox{}\fi\par}{}{}
\makeatother
% Allow footnotes in longtable head/foot
\IfFileExists{footnotehyper.sty}{\usepackage{footnotehyper}}{\usepackage{footnote}}
\makesavenoteenv{longtable}

\usepackage{graphicx}
\makeatletter
\newsavebox\pandoc@box
\newcommand*\pandocbounded[1]{% scales image to fit in text height/width
  \sbox\pandoc@box{#1}%
  \Gscale@div\@tempa{\textheight}{\dimexpr\ht\pandoc@box+\dp\pandoc@box\relax}%
  \Gscale@div\@tempb{\linewidth}{\wd\pandoc@box}%
  \ifdim\@tempb\p@<\@tempa\p@\let\@tempa\@tempb\fi% select the smaller of both
  \ifdim\@tempa\p@<\p@\scalebox{\@tempa}{\usebox\pandoc@box}%
  \else\usebox{\pandoc@box}%
  \fi%
}
% Set default figure placement to htbp
\def\fps@figure{htbp}
\makeatother


% definitions for citeproc citations
\NewDocumentCommand\citeproctext{}{}
\NewDocumentCommand\citeproc{mm}{%
  \begingroup\def\citeproctext{#2}\cite{#1}\endgroup}
\makeatletter
 % allow citations to break across lines
 \let\@cite@ofmt\@firstofone
 % avoid brackets around text for \cite:
 \def\@biblabel#1{}
 \def\@cite#1#2{{#1\if@tempswa , #2\fi}}
\makeatother
\newlength{\cslhangindent}
\setlength{\cslhangindent}{1.5em}
\newlength{\csllabelwidth}
\setlength{\csllabelwidth}{3em}
\newenvironment{CSLReferences}[2] % #1 hanging-indent, #2 entry-spacing
 {\begin{list}{}{%
  \setlength{\itemindent}{0pt}
  \setlength{\leftmargin}{0pt}
  \setlength{\parsep}{0pt}
  % turn on hanging indent if param 1 is 1
  \ifodd #1
   \setlength{\leftmargin}{\cslhangindent}
   \setlength{\itemindent}{-1\cslhangindent}
  \fi
  % set entry spacing
  \setlength{\itemsep}{#2\baselineskip}}}
 {\end{list}}
\usepackage{calc}
\newcommand{\CSLBlock}[1]{\hfill\break\parbox[t]{\linewidth}{\strut\ignorespaces#1\strut}}
\newcommand{\CSLLeftMargin}[1]{\parbox[t]{\csllabelwidth}{\strut#1\strut}}
\newcommand{\CSLRightInline}[1]{\parbox[t]{\linewidth - \csllabelwidth}{\strut#1\strut}}
\newcommand{\CSLIndent}[1]{\hspace{\cslhangindent}#1}


\usepackage[nolongtablepatch]{lineno}
\linenumbers



\usepackage{newtx}

\defaultfontfeatures{Scale=MatchLowercase}
\defaultfontfeatures[\rmfamily]{Ligatures=TeX,Scale=1}





\title{eyetools: an R package for open-source analysis of eye data}


\shorttitle{eyetools: eye data analysis}


\usepackage{etoolbox}








\authorsnames{Tom Beesley,Matthew Ivory}





\affiliation{
{Lancaster University}}




\leftheader{Beesley and Ivory}



\abstract{Eye data analysis has become an integral part of research in
both academic and commercial settings. In the behavioural sciences there
is an ever increasing desire and need to shift analysis from proprietary
software to open source alternatives. Here we introduce eyetools, an R
package that provides an open source means of conducting eye data
analysis. eyetools is aimed at researchers who may not have experience
programming their own eye data analysis routines. The package provides a
number of functions that facilitate key steps in the pipeline of eye
data analysis, from basic data manipulation, extraction of fixations and
saccades, to trial level summaries and visualisations. In this article
we introduce these core features of the package and provide a
step-by-step guide which will enable researchers to find open source
solutions to their eye data analysis. We discuss the ways in which we
believe eyetools might be expanded in the future and how it may provide
a platform for collaborative work on eye data analysis.}

\keywords{software; R; eye-tracking; fixations; saccades;
areas-of-interest}

\authornote{\par{\addORCIDlink{Tom Beesley}{0000-0003-2836-2743}} 

\par{       }
\par{Correspondence concerning this article should be addressed to Tom
Beesley, Lancaster University, Department of Psychology, Lancaster
University, UK, LA1 4YD, UK, Email: t.beesley@lancaster.ac.uk}
}

\makeatletter
\let\endoldlt\endlongtable
\def\endlongtable{
\hline
\endoldlt
}
\makeatother

\urlstyle{same}



\makeatletter
\@ifpackageloaded{caption}{}{\usepackage{caption}}
\AtBeginDocument{%
\ifdefined\contentsname
  \renewcommand*\contentsname{Table of contents}
\else
  \newcommand\contentsname{Table of contents}
\fi
\ifdefined\listfigurename
  \renewcommand*\listfigurename{List of Figures}
\else
  \newcommand\listfigurename{List of Figures}
\fi
\ifdefined\listtablename
  \renewcommand*\listtablename{List of Tables}
\else
  \newcommand\listtablename{List of Tables}
\fi
\ifdefined\figurename
  \renewcommand*\figurename{Figure}
\else
  \newcommand\figurename{Figure}
\fi
\ifdefined\tablename
  \renewcommand*\tablename{Table}
\else
  \newcommand\tablename{Table}
\fi
}
\@ifpackageloaded{float}{}{\usepackage{float}}
\floatstyle{ruled}
\@ifundefined{c@chapter}{\newfloat{codelisting}{h}{lop}}{\newfloat{codelisting}{h}{lop}[chapter]}
\floatname{codelisting}{Listing}
\newcommand*\listoflistings{\listof{codelisting}{List of Listings}}
\makeatother
\makeatletter
\makeatother
\makeatletter
\@ifpackageloaded{caption}{}{\usepackage{caption}}
\@ifpackageloaded{subcaption}{}{\usepackage{subcaption}}
\makeatother
\makeatletter
\@ifpackageloaded{fontawesome5}{}{\usepackage{fontawesome5}}
\makeatother

% From https://tex.stackexchange.com/a/645996/211326
%%% apa7 doesn't want to add appendix section titles in the toc
%%% let's make it do it
\makeatletter
\xpatchcmd{\appendix}
  {\par}
  {\addcontentsline{toc}{section}{\@currentlabelname}\par}
  {}{}
\makeatother

%% Disable longtable counter
%% https://tex.stackexchange.com/a/248395/211326

\usepackage{etoolbox}

\makeatletter
\patchcmd{\LT@caption}
  {\bgroup}
  {\bgroup\global\LTpatch@captiontrue}
  {}{}
\patchcmd{\longtable}
  {\par}
  {\par\global\LTpatch@captionfalse}
  {}{}
\apptocmd{\endlongtable}
  {\ifLTpatch@caption\else\addtocounter{table}{-1}\fi}
  {}{}
\newif\ifLTpatch@caption
\makeatother

\begin{document}

\maketitle


\setcounter{secnumdepth}{-\maxdimen} % remove section numbering

\setlength\LTleft{0pt}

\resetlinenumber[1]

Eye-tracking is now an established and widely used tool that provides
powerful measurements of human behaviour. Its application across
academic research is far-reaching, with major importance to the fields
of computer science and AI, economics, psychology, and even having
influence on art and design. Beyond purely academic application,
eye-tracking is widely used in commercial fields where it can provide
insights into consumer behaviour, helping to shape product design and
marketing.

By recording the movement of an individual's gaze during research
studies, users can quantify where and how long individual's look at
particular regions of space (usually with a focus on stimuli presented
on a 2D screen, but also within 3D space). Eye-tracking provides a rich
stream of continuous data which can offer powerful insights into
real-time cognitive processing. Undoubtedly, Psychology (and related
broad areas of research such as vision science, linguistics, etc) is
perhaps the scientific discipline which has seen the most substantial
adoption of eye-data research, where eye-tracking systems are now
commonplace in centres of academic research. Such data allow researchers
to inspect the interplay of cognitive processes such as attention,
memory, decision making, and reading, with high temporal precision
(\citeproc{ref-beesley2019}{Beesley et al., 2019}).

While there are abundant uses and benefits of collecting eye-movement
data, the collection of such data has many implications for the eventual
analysis work that is undertaken by the researcher. The continual stream
of recording can lead to an overwhelming amount of raw data: modern
eye-trackers can record data at 1000 Hz and above, which results in 3.6
million samples taken per hour. The continual nature of the data also
leads to a wealth of choice in the manner in which it is analysed. As
such, the provision of suitable computational software for data
reduction and processing is an important part of eye-tracking research.
The companies behind eye-tracking devices offer licensed software that
will perform many of the common steps for eye-data analysis. However,
there are several disadvantages to using such proprietary software in a
research context. Firstly, the software will typically have an ongoing
license cost for continual use. Secondly, the algorithms driving the
operations within such software are not readily available for inspection
or adaptation for new purposes. These significant constraints mean that
the use of proprietary analysis software will lead to a failure to meet
the basic open-science principle of analysis reproduction, for example
as set out by the UK Reproducibility Network: ``We expect researchers
to\ldots{} make their research methods, software, outputs and data open,
and available at the earliest possible point\ldots The reproducibility
of both research methods and research results \ldots is critical to
research in certain contexts, particularly in the experimental sciences
with a quantitative focus\ldots{}''.

In the current article we introduce a new toolkit for eye-data
processing and analysis called ``\emph{eyetools}'', which takes the form
of an R package. R packages (like R itself) are free to use and are
therefore available for users across the world. The package provides a
(growing) number of functions that provide an efficient and effective
means to conduct basic eye-data analysis. \emph{eyetools} is built with
academic researchers in the psychological sciences in mind, though there
is no reason why the package would not be effective more generally. The
functions within the package reflect steps in a comprehensive analysis
pipeline, taking the user from initial handling of raw eye data, to
summarising data for each period of a procedure, to the visualisation of
the data in plots. Since the pipeline is contained within the one
package, it is not reliant on external software, which enables easy
reproducibility of any analyses. Importantly, the functions are simple
to use, ensuring that the package will be beneficial for researchers who
are unfamiliar with eye data analysis. It should also appeal to
researchers accustomed to working with eye data in other environments
who wish to transfer to working in R.

\emph{eyetools} is, of course, not the only package in R that allows
users to work with eye data. We conducted a survey of the available
packages on CRAN (The Comprehensive R Archive Network), by searching the
list of available packages
(\url{https://cran.r-project.org/web/packages/available_packages_by_name.html})
for the terms ``eye'' (which covers terms ``eye-tracking'' and
``eyetracking''), ``gaze'', ``saccades'' and ``fixation''. We considered
packages that were designed for work within the psychological
sciences\footnote{We did not present or summarise here packages that
  were designed for work in opthalmology research.}, which identified 12
packages that offer relevant functions for the analysis of eye data,
shown in Table 1. \textbf{\emph{eyeTrackr}}, \textbf{\emph{eyelinker}},
and \textbf{\emph{eyelinkReader}}, all offer functionality for data only
from experiments that have used `EyeLink' trackers (S-R Research). In
contrast, \textbf{\emph{eyetools}} provides functions that are
hardware-agnostic, relying on a format of data that can be achieved from
any source. The \textbf{\emph{emov}} package offers a limited set of
functions and is primarily designed for fixation detection, using the
same dispersion method employed in \textbf{\emph{eyetools}}.
\textbf{\emph{eyetrackingR}} and \textbf{\emph{kollaR}} are perhaps the
most comprehensive alternative packages and having some similarity to
eyetools in their broad use-case. \textbf{\emph{eyetrackingR}} offers a
large suite of functionality and, like \textbf{\emph{eyetools}}, can be
applied across the entire pipeline. It has functions for cleaning data
and various plotting functions, including analysis over time. It does
not feature algorithms regarding the detection of events such as
saccades or fixations. This limits the ability to conduct bespoke
analysis steps and it means that analysis needs to be conducted on raw
data. This is disadvantageous both in terms of computing time and in the
open sharing of data (event data are an order of magnitude smaller in
size than raw data). kollaR is hardware-agnostic and offers functions
for data processing, fixation detection and some plotting features.
Several packages offer specific use-cases: \textbf{\emph{saccadr}} is a
package that can be used to extract saccades from raw data, and has
functionality to convert binocular data into monocular data;
\textbf{\emph{fpa}} offers a particular spatio-temporal analysis and
associated plots; \textbf{\emph{gazepath}} offers detection of fixations
and saccades, in addition to an optional GUI. The packages
\textbf{\emph{eyeris}} and \textbf{\emph{Pupillometry}} are focussed on
the analysis of pupillometry data, which is not a focus of the eyetools
package (at present).

{\def\LTcaptype{none} % do not increment counter
\begin{longtable}[]{@{}
  >{\raggedright\arraybackslash}p{(\linewidth - 16\tabcolsep) * \real{0.1111}}
  >{\centering\arraybackslash}p{(\linewidth - 16\tabcolsep) * \real{0.1111}}
  >{\centering\arraybackslash}p{(\linewidth - 16\tabcolsep) * \real{0.1111}}
  >{\centering\arraybackslash}p{(\linewidth - 16\tabcolsep) * \real{0.1111}}
  >{\centering\arraybackslash}p{(\linewidth - 16\tabcolsep) * \real{0.1111}}
  >{\centering\arraybackslash}p{(\linewidth - 16\tabcolsep) * \real{0.1111}}
  >{\centering\arraybackslash}p{(\linewidth - 16\tabcolsep) * \real{0.1111}}
  >{\centering\arraybackslash}p{(\linewidth - 16\tabcolsep) * \real{0.1111}}
  >{\centering\arraybackslash}p{(\linewidth - 16\tabcolsep) * \real{0.1111}}@{}}
\toprule\noalign{}
\begin{minipage}[b]{\linewidth}\raggedright
Package
\end{minipage} & \begin{minipage}[b]{\linewidth}\centering
Hardware-agnostic
\end{minipage} & \begin{minipage}[b]{\linewidth}\centering
Data Import
\end{minipage} & \begin{minipage}[b]{\linewidth}\centering
Data processing
\end{minipage} & \begin{minipage}[b]{\linewidth}\centering
Fixation detection
\end{minipage} & \begin{minipage}[b]{\linewidth}\centering
Saccade detection
\end{minipage} & \begin{minipage}[b]{\linewidth}\centering
Pupillometry
\end{minipage} & \begin{minipage}[b]{\linewidth}\centering
Plotting
\end{minipage} & \begin{minipage}[b]{\linewidth}\centering
Inferential Analysis
\end{minipage} \\
\midrule\noalign{}
\endhead
\bottomrule\noalign{}
\endlastfoot
emov & \faIcon{check} & & \faIcon{check} & \faIcon{check} &
\faIcon{check} & & & \\
eyelinker & & Eyelink & & & & & & \\
eyelinkReader & & EyeLink & & \faIcon{check} & \faIcon{check} & &
\faIcon{check} & \\
eyeris & \faIcon{check} & EyeLink & \faIcon{check} & & & \faIcon{check}
& \faIcon{check} & \\
eyetools & \faIcon{check} & Tobii & \faIcon{check} & \faIcon{check} &
\faIcon{check} & & \faIcon{check} & \\
eyetrackingR & \faIcon{check} & & \faIcon{check} & & & & \faIcon{check}
& \faIcon{check} \\
eyeTrackR & & EyeLink & \faIcon{check} & & & & & \\
fpa & \faIcon{check} & & \faIcon{check} & & & & \faIcon{check} & \\
gazepath & \faIcon{check} & & & \faIcon{check} & & & & \\
kollaR & \faIcon{check} & & \faIcon{check} & \faIcon{check} & & &
\faIcon{check} & \\
PuppilometryR & & & \faIcon{check} & & & \faIcon{check} & \faIcon{check}
& \\
saccadr & \faIcon{check} & & & \faIcon{check} & & & & \\
\end{longtable}
}

\protect\phantomsection\label{refs}
\begin{CSLReferences}{1}{0}
\bibitem[\citeproctext]{ref-beesley2019}
Beesley, T., Pearson, D., \& Le Pelley, M. (2019). \emph{Chapter 1 - eye
tracking as a tool for examining cognitive processes} (G. Foster, Ed.;
pp. 1--30). Academic Press.
\url{https://doi.org/10.1016/B978-0-12-813092-6.00002-2}

\end{CSLReferences}






\end{document}
